% Options for packages loaded elsewhere
\PassOptionsToPackage{unicode}{hyperref}
\PassOptionsToPackage{hyphens}{url}
\PassOptionsToPackage{dvipsnames,svgnames,x11names}{xcolor}
\documentclass[
  british,
  11pt,
]{article}
\usepackage{xcolor}
\usepackage[margin=25mm]{geometry}
\usepackage{amsmath,amssymb}
\setcounter{secnumdepth}{-\maxdimen} % remove section numbering
\usepackage{iftex}
\ifPDFTeX
  \usepackage[T1]{fontenc}
  \usepackage[utf8]{inputenc}
  \usepackage{textcomp} % provide euro and other symbols
\else % if luatex or xetex
  \usepackage{unicode-math} % this also loads fontspec
  \defaultfontfeatures{Scale=MatchLowercase}
  \defaultfontfeatures[\rmfamily]{Ligatures=TeX,Scale=1}
\fi
\usepackage{lmodern}
\ifPDFTeX\else
  % xetex/luatex font selection
  \setmainfont[]{TeX Gyre Termes}
  \setmathfont[]{TeX Gyre Termes Math}
\fi
% Use upquote if available, for straight quotes in verbatim environments
\IfFileExists{upquote.sty}{\usepackage{upquote}}{}
\IfFileExists{microtype.sty}{% use microtype if available
  \usepackage[]{microtype}
  \UseMicrotypeSet[protrusion]{basicmath} % disable protrusion for tt fonts
}{}
\makeatletter
\@ifundefined{KOMAClassName}{% if non-KOMA class
  \IfFileExists{parskip.sty}{%
    \usepackage{parskip}
  }{% else
    \setlength{\parindent}{0pt}
    \setlength{\parskip}{6pt plus 2pt minus 1pt}}
}{% if KOMA class
  \KOMAoptions{parskip=half}}
\makeatother
\usepackage{longtable,booktabs,array}
\newcounter{none} % for unnumbered tables
\usepackage{calc} % for calculating minipage widths
% Correct order of tables after \paragraph or \subparagraph
\usepackage{etoolbox}
\makeatletter
\patchcmd\longtable{\par}{\if@noskipsec\mbox{}\fi\par}{}{}
\makeatother
% Allow footnotes in longtable head/foot
\IfFileExists{footnotehyper.sty}{\usepackage{footnotehyper}}{\usepackage{footnote}}
\makesavenoteenv{longtable}
% definitions for citeproc citations
\NewDocumentCommand\citeproctext{}{}
\NewDocumentCommand\citeproc{mm}{%
  \begingroup\def\citeproctext{#2}\cite{#1}\endgroup}
\makeatletter
 % allow citations to break across lines
 \let\@cite@ofmt\@firstofone
 % avoid brackets around text for \cite:
 \def\@biblabel#1{}
 \def\@cite#1#2{{#1\if@tempswa , #2\fi}}
\makeatother
\newlength{\cslhangindent}
\setlength{\cslhangindent}{1.5em}
\newlength{\csllabelwidth}
\setlength{\csllabelwidth}{3em}
\newenvironment{CSLReferences}[2] % #1 hanging-indent, #2 entry-spacing
 {\begin{list}{}{%
  \setlength{\itemindent}{0pt}
  \setlength{\leftmargin}{0pt}
  \setlength{\parsep}{0pt}
  % turn on hanging indent if param 1 is 1
  \ifodd #1
   \setlength{\leftmargin}{\cslhangindent}
   \setlength{\itemindent}{-1\cslhangindent}
  \fi
  % set entry spacing
  \setlength{\itemsep}{#2\baselineskip}}}
 {\end{list}}
\usepackage{calc}
\newcommand{\CSLBlock}[1]{\hfill\break\parbox[t]{\linewidth}{\strut\ignorespaces#1\strut}}
\newcommand{\CSLLeftMargin}[1]{\parbox[t]{\csllabelwidth}{\strut#1\strut}}
\newcommand{\CSLRightInline}[1]{\parbox[t]{\linewidth - \csllabelwidth}{\strut#1\strut}}
\newcommand{\CSLIndent}[1]{\hspace{\cslhangindent}#1}
\ifLuaTeX
\usepackage[bidi=basic,shorthands=off]{babel}
\else
\usepackage[bidi=default,shorthands=off]{babel}
\fi
\ifPDFTeX
\else
\babelfont{rm}[]{TeX Gyre Termes}
\fi
\ifLuaTeX
  \usepackage{selnolig} % disable illegal ligatures
\fi
\setlength{\emergencystretch}{3em} % prevent overfull lines
\providecommand{\tightlist}{%
  \setlength{\itemsep}{0pt}\setlength{\parskip}{0pt}}
\usepackage{microtype}
\usepackage{needspace}
\setlength{\emergencystretch}{3em}
\widowpenalty=10000
\clubpenalty=10000
\usepackage{bookmark}
\IfFileExists{xurl.sty}{\usepackage{xurl}}{} % add URL line breaks if available
\urlstyle{same}
\hypersetup{
  pdftitle={Supporting Information: Incremental Telemetry for Wind-Power Shortfall Decisions under a Fixed Forecast Commitment},
  pdflang={en-GB},
  colorlinks=true,
  linkcolor={Maroon},
  filecolor={Maroon},
  citecolor={Blue},
  urlcolor={Blue},
  pdfcreator={LaTeX via pandoc}}

\title{Supporting Information: Incremental Telemetry for Wind-Power
Shortfall Decisions under a Fixed Forecast Commitment}
\author{}
\date{}

\begin{document}
\maketitle

\section{S1. Shared information and registered
comparisons}\label{s1.-shared-information-and-registered-comparisons}

\subsection{S1.1. A common forecast, followed by an incremental decision
layer}\label{s1.1.-a-common-forecast-followed-by-an-incremental-decision-layer}

The study evaluates additional telemetry at a reserve-decision layer,
conditional on one corresponding unchanged upstream power forecast. Let
\(t\) denote a decision issue on the ten-minute grid. The upstream
forecast uses 36 wind-speed and active-power observations ending at
\(t-6\) and predicts power at \(t+6\). Thus, the reserve decision has a
six-step lead, whereas the common forecast has a twelve-step lead from
its last input. The target is one ten-minute power interval, not energy
delivered throughout the following hour. The access contract supplies
this single forecast rather than the complete set of potentially arrived
forecasts, errors or controller messages. It is a controlled design
comparison, not a reconstruction of an operator's complete information
set.

Every direct decision estimator receives four common inputs: the shared
forecast, the sine and cosine of issue time within the day, and
normalised turbine identity. We denote this information by
\textbf{Common}. The name \textbf{P1} means that current active power is
added to Common at the decision layer. It does not mean that the
complete forecasting and decision pipeline uses active power alone. In
particular, the shared forecast still requires historical wind speed and
active power.

The primary condition withholds pitch from the decision layer and
retains fresh wind speed, active power and reactive power as candidate
inputs. The complete registered lattice consists of Common and each of
the seven non-empty combinations of wind speed (\(W\)), active power
(\(P\)) and reactive power (\(Q\)), with 1, 6 or 36 last-arrived
observations. Each recipe is fitted with histogram gradient boosting
(HGB) and a multilayer perceptron (MLP). Table S1 lists all 22 input
recipes before the estimator and calibration dimensions are added.

\textbf{Table S1. Complete pitch-withheld input lattice.} Width is the
number of raw inputs selected by the reserve learner before its
missing-value transformation. All rows share the forecast, clock terms
and turbine identity. The entries in the last three columns correspond
to lookbacks of 1, 6 and 36 observations, respectively.

{\def\LTcaptype{none} % do not increment counter
\begin{longtable}[]{@{}
  >{\raggedright\arraybackslash}p{(\linewidth - 10\tabcolsep) * \real{0.1667}}
  >{\raggedright\arraybackslash}p{(\linewidth - 10\tabcolsep) * \real{0.1667}}
  >{\raggedleft\arraybackslash}p{(\linewidth - 10\tabcolsep) * \real{0.1667}}
  >{\raggedleft\arraybackslash}p{(\linewidth - 10\tabcolsep) * \real{0.1667}}
  >{\raggedleft\arraybackslash}p{(\linewidth - 10\tabcolsep) * \real{0.1667}}
  >{\raggedleft\arraybackslash}p{(\linewidth - 10\tabcolsep) * \real{0.1667}}@{}}
\toprule\noalign{}
\begin{minipage}[b]{\linewidth}\raggedright
Added channels
\end{minipage} & \begin{minipage}[b]{\linewidth}\raggedright
Registered labels
\end{minipage} & \begin{minipage}[b]{\linewidth}\raggedleft
Incremental channel count
\end{minipage} & \begin{minipage}[b]{\linewidth}\raggedleft
Width, 1 observation
\end{minipage} & \begin{minipage}[b]{\linewidth}\raggedleft
Width, 6 observations
\end{minipage} & \begin{minipage}[b]{\linewidth}\raggedleft
Width, 36 observations
\end{minipage} \\
\midrule\noalign{}
\endhead
\bottomrule\noalign{}
\endlastfoot
None & Common & 0 & 4 & --- & --- \\
W & W1, W6, W36 & 1 & 5 & 10 & 40 \\
P & P1, P6, P36 & 1 & 5 & 10 & 40 \\
Q & Q1, Q6, Q36 & 1 & 5 & 10 & 40 \\
W, P & WP1, WP6, WP36 & 2 & 6 & 16 & 76 \\
W, Q & WQ1, WQ6, WQ36 & 2 & 6 & 16 & 76 \\
P, Q & PQ1, PQ6, PQ36 & 2 & 6 & 16 & 76 \\
W, P, Q & WPQ1, WPQ6, WPQ36 & 3 & 7 & 22 & 112 \\
\end{longtable}
}

Each fitted recipe has five historical calibration offsets, indexed by
\(q\in\{0.90,0.925,0.95,0.975,0.99\}\). A calibrated persistence
comparator is also retained. The pitch-withheld full family therefore
contains \((22\times2+1)\times5=225\) policies. The five archived
learning instances are evaluated for each policy; they are not
additional policies in this count.

\subsection{S1.2. Secondary information conditions and
persistence}\label{s1.2.-secondary-information-conditions-and-persistence}

The five pitch-observed secondary conditions use a registered partial
lattice: Common, P1, WP1, WP6, WP36 and WPQ36, each with HGB and MLP,
plus persistence. With the same five offset levels, each condition
contains 65 policies. The six conditions together contain 550 policy
labels in the full-q evaluation. Their labels, stored indices and
resource tuples are provided in
\texttt{supplement\_data/registered\_confirmation\_policies.csv}.

\textbf{Table S2. Registered information conditions.} The shifts apply
to reserve-layer telemetry. A fresh upstream feed is assumed for the
shared forecast in every condition. Delays and withholding are imposed
evaluation conditions; their frequency in operational wind fleets is not
measured here.

{\def\LTcaptype{none} % do not increment counter
\begin{longtable}[]{@{}
  >{\raggedright\arraybackslash}p{(\linewidth - 6\tabcolsep) * \real{0.2500}}
  >{\raggedleft\arraybackslash}p{(\linewidth - 6\tabcolsep) * \real{0.2500}}
  >{\raggedright\arraybackslash}p{(\linewidth - 6\tabcolsep) * \real{0.2500}}
  >{\raggedright\arraybackslash}p{(\linewidth - 6\tabcolsep) * \real{0.2500}}@{}}
\toprule\noalign{}
\begin{minipage}[b]{\linewidth}\raggedright
Condition
\end{minipage} & \begin{minipage}[b]{\linewidth}\raggedleft
Reserve-feed delay
\end{minipage} & \begin{minipage}[b]{\linewidth}\raggedright
Pitch history in direct estimators
\end{minipage} & \begin{minipage}[b]{\linewidth}\raggedright
Registered direct-input recipes
\end{minipage} \\
\midrule\noalign{}
\endhead
\bottomrule\noalign{}
\endlastfoot
Pitch withheld & 0 min & No & All 22 recipes in Table S1 \\
Fresh & 0 min & Yes & Common, P1, WP1, WP6, WP36, WPQ36 \\
Delay 1 & 10 min & Yes & Same six-recipe partial lattice \\
Delay 2 & 20 min & Yes & Same six-recipe partial lattice \\
Delay 3 & 30 min & Yes & Same six-recipe partial lattice \\
Delay 6 stress & 60 min & Yes & Same six-recipe partial lattice \\
\end{longtable}
}

The event-time window is always \([t-35,t]\). For delay \(d\) and
requested lookback \(L\), retained observations end at \(t-d\) and begin
no earlier than \(t-35\). Unavailable observations are not recovered by
shifting the complete history window backwards. When pitch is permitted,
all legally available pitch observations within this same window are
shared by every direct candidate. Consequently, a label such as P1 in a
pitch-observed condition includes pitch history in addition to its named
input; it cannot be interpreted as the same complete information set as
pitch-withheld P1.

The persistence raw action is
\(\max(F_{i,t}-P^{\mathrm{last}}_{i,t},0)\), with zero used when the
last active-power observation is missing, followed by the same
historical offset family. This comparator uses the shared forecast and
last-arrived power, and does not use the clock, turbine identity or
available pitch. It is a simple functional comparator rather than an
equal-input estimator control. Its five stored copies implement the same
deterministic function and do not supply five independent replications.

Resource order is lexicographic in (incremental telemetry channels,
learned predictor passes, raw input width). The persistence tuple is
\((1,0,2)\); a pitch-withheld P1 direct estimator has \((1,1,5)\);
Common has \((0,1,4)\). HGB and MLP with identical selected inputs are
tied under this declared order. The order does not rank parameter count,
inference time, communication traffic or total system cost. The
implementation constructs the complete feature matrix before selecting
columns, so functional input sparsity is not evidence of an implemented
end-to-end speedup.

\section{S2. Data, chronology and policy
construction}\label{s2.-data-chronology-and-policy-construction}

\subsection{S2.1. Coordinate and population
rules}\label{s2.1.-coordinate-and-population-rules}

The Hill of Towie analysis uses the registered 21 turbines and the 2023,
2024 and 2025 annual SCADA archives\textsuperscript{1}. Records are
placed on a UTC ten-minute period-end grid. Wind speed comes from
\texttt{wtc\_AcWindSp\_mean}, active and reactive power from
\texttt{wtc\_ActPower\_mean} and \texttt{wtc\_ReactPwr\_mean}, and pitch
is the mean of \texttt{wtc\_PitcPosA\_mean},
\texttt{wtc\_PitcPosB\_mean} and \texttt{wtc\_PitcPosC\_mean} when all
three blade values are finite. Negative wind speed is marked missing;
finite negative active or reactive power is retained. Missing
time--turbine coordinates remain missing rather than being filled with
invented measurements.

Exact redundant source records are accepted only after equality of the
original six consumed fields, including matching non-finite values, has
been checked. Any conflicting consumed value rejects the source run.
Equality at the cross-year boundary is checked before wind-speed masking
or blade-pitch averaging. The earlier failure of a stricter
raw-row-uniqueness check is retained in the source audit; accepting
equal copies does not modify the annual archives.

All compared policies use the same anchors, shared forecast, shortfall
target and target mask. Eligibility requires a finite future
active-power observation; it does not require finite pitch or finite
inputs to any particular estimator. Zero and negative power targets
remain eligible. No alarm, operating-state, control-intervention or
policy-performance filter is applied. The 2025 confirmation population
contains 1,059,248 valid turbine intervals per archived instance.
Inference is conditional on this recorded-target population and does not
cover unrecorded future power.

WTB supplies a separate development study with 134 turbines and 245
index days from the SDWPF dataset\textsuperscript{2}. Its evaluation
period had been inspected before the present explanation and is not an
untouched confirmation set. The same principle retains finite zero and
negative future power and uses a common target population across input
recipes. WTB and Hill of Towie are not pooled into a single effect
estimate.

\Needspace{18\baselineskip}

\subsection{S2.2. Learning, calibration and selection
windows}\label{s2.2.-learning-calibration-and-selection-windows}

\textbf{Table S3. Registered stages.} Calendar windows are half-open on
the issue grid. The final issues are removed until their target plus one
release step precedes the next stage. Historical observations before a
stage may be used when they have already arrived.

{\def\LTcaptype{none} % do not increment counter
\begin{longtable}[]{@{}
  >{\raggedright\arraybackslash}p{(\linewidth - 4\tabcolsep) * \real{0.3333}}
  >{\raggedright\arraybackslash}p{(\linewidth - 4\tabcolsep) * \real{0.3333}}
  >{\raggedright\arraybackslash}p{(\linewidth - 4\tabcolsep) * \real{0.3333}}@{}}
\toprule\noalign{}
\begin{minipage}[b]{\linewidth}\raggedright
Stage
\end{minipage} & \begin{minipage}[b]{\linewidth}\raggedright
Hill of Towie
\end{minipage} & \begin{minipage}[b]{\linewidth}\raggedright
WTB
\end{minipage} \\
\midrule\noalign{}
\endhead
\bottomrule\noalign{}
\endlastfoot
Shared point-forecast fit & 2023-01-01 to 2023-07-01 & Days 1--90 \\
Reserve-estimator fit & 2023-07-01 to 2024-01-01 & Days 136--175 \\
Offset fit & 2024-01-01 to 2024-07-01 & Days 176--192 \\
Historical q and recipe selection & 2024-07-01 to 2025-01-01 & Days
193--210 \\
Evaluation & 2025-01-01 to 2026-01-01, sealed confirmation & Days
211--245, development \\
\end{longtable}
}

The shared point forecaster is a squared-error HGB using only wind-speed
and active-power histories. Hill of Towie uses 120,000 historical
point-fit rows and seed 20261018. Because its inputs end at forecast
publication, the fitted target is twelve steps after that input anchor.
The last point-fit label is required to mature before the first forecast
publication used for policy fitting. This additional boundary check
prevents a forecast model from using labels that would not yet have
existed at its first use.

Hill of Towie reserve models use 360,000 finite-target fit rows per
instance. The registered instance seeds are 201--205. For each seed,
every input subset and estimator uses the same sampled rows. WTB uses
nested prefixes of 10,000, 40,000, 120,000 and 360,000 rows for pitch
withholding; the five pitch-observed conditions use only 120,000 and
360,000 rows. These are 14 condition/budget settings, not a complete
condition-by-four-budget factorial. Increasing row count changes both
data density and training computation within an unchanged calendar; it
does not add evaluation weather days.

\subsection{S2.3. Estimation and missing
inputs}\label{s2.3.-estimation-and-missing-inputs}

HGB and MLP both fit the shortfall's 0.90 conditional quantile using
pinball loss. HGB uses 100 boosting iterations, at most 15 leaves,
minimum leaf size 50, learning rate 0.08 and L2 regularisation 1.0, with
early stopping disabled. Its historical pipeline imputes missing inputs
with zero and adds missing-value indicators, retaining structurally
empty features. The imputer is fitted only on the corresponding
historical fit rows.

The MLP has two ReLU hidden layers of widths 64 and 32 and a linear
output. It uses 80 fixed epochs, batch size 2048 and AdamW with learning
rate 0.001 and weight decay 0.0001. Input means and scales use only
finite entries in the actual fit rows; missing standardised inputs are
set to zero, with an indicator retained for every raw feature. Target
scale is the fit-target standard deviation, floored at one. The final
bias is initialised using the fit-target 0.90 quantile. Training uses
deterministic CPU settings. Neither estimator is tuned against 2025.

Table S1 reports raw selected width before these missing-value
transformations. HGB's added indicator count depends on historical
missingness, whereas the MLP retains a full set of indicators. The raw
width therefore is not a claim about identical transformed dimensions or
equal computational cost across estimators.

\subsection{S2.4. Historical calibration and
selection}\label{s2.4.-historical-calibration-and-selection}

For a fitted raw prediction \(a_{i,t}\), let
\(a^+_{i,t}=\max(a_{i,t},0)\). Each historical offset is the higher
sample quantile

\[
b_q=Q_q^{\mathrm{higher}}\{S_{i,t}-a^+_{i,t}: (i,t)\text{ in the offset-fit population}\},
\qquad R_{i,t,q}=\max(a^+_{i,t}+b_q,0).
\]

Offsets may be negative. The grid \(q=0.90,0.925,0.95,0.975,0.99\)
changes the historical offset, not the quantile used to train five
different networks. Thus P1/HGB/q0.95 and P1/HGB/q0.925 share the same
fitted model and have different precomputed historical offsets. Neither
is fitted again after confirmation is opened.

Historical q selection uses one q shared by all five instances of a
recipe. Eligible levels must satisfy the historical simultaneous
event-rate upper limit of 0.10 for every instance. The eligible level
with the lowest five-instance mean historical cost is selected; ties use
the lower grid index. If no level qualifies, q0.99 is retained as an
explicitly unqualified diagnostic rather than a successful choice.

The historical recipe-family assessment then applies the risk and
near-optimality rules in Section S3 to those fixed q choices, retaining
all minimum-resource ties or an empty set. The reference is the
lowest-cost historically qualified WPQ36 HGB or MLP recipe among its
historically selected q levels, with deterministic label tie-breaking.
If no such reference qualifies, the registered q0.99 fallback remains
unqualified and cannot establish primary confirmation. In the completed
pitch-withheld selection, the frozen candidate was P1/HGB/q0.95 and the
frozen reference was WPQ36/HGB/q0.95.

\section{S3. Inference, confirmation and finite-family
identification}\label{s3.-inference-confirmation-and-finite-family-identification}

\subsection{S3.1. Quantities and approximate joint
bounds}\label{s3.1.-quantities-and-approximate-joint-bounds}

For the shared forecast \(F_{i,t}\) and future power \(P_{i,t+6}\),
shortfall is \(S_{i,t}=\max(F_{i,t}-P_{i,t+6},0)\). Define uncovered
shortfall \(U_{i,t}=(S_{i,t}-R_{i,t})_+\) and event indicator
\(1(S_{i,t}>R_{i,t})\). The loss for one valid turbine interval is

\[
\ell_{i,t}=\frac{R_{i,t}+10U_{i,t}}{6}.
\]

The factor \(1/6\) converts ten-minute power quantities to
interval-equivalent energy. These are turbine-level procurement and
shortage proxies, not measured reserve purchases, cash costs or net
shortfall at a point of common coupling. Recipe cost \(C_m\) is the mean
interval loss, averaged over the five fixed archived fits. Event rates
\(V_{s,m}\) are checked separately for each fit \(s\). The cost of
averaged actions is not evaluated: averaging five realised losses does
not generally equal the loss of an averaged reserve action. No arbitrary
future training instance or ensemble deployment is certified.

To distinguish a finite-record summary from a population interpretation,
let \(N_d\) be the number of eligible turbine targets on UTC day \(d\),
\(A_{s,m,d}\) the stored loss total, and \(B_{s,m,d}\) the stored
strict-event count. With \(K=5\) fixed fits, the observed quantities are

\[
\widehat c_m=\frac{\sum_d K^{-1}\sum_{s=1}^K A_{s,m,d}}{\sum_d N_d},
\qquad
\widehat v_{s,m}=\frac{\sum_d B_{s,m,d}}{\sum_d N_d}.
\]

These are the reported \(C_m\) and \(V_{s,m}\); they are deterministic
once the finite recorded targets and actions are given. The bootstrap
does not create uncertainty about their arithmetic. A working process
interpretation instead treats the daily vectors
\((N_d,\{A_{s,m,d},B_{s,m,d}\}_{s,m})\) as observations from a
repeatable, weakly dependent daily-sequence law \(\mathcal D\),
conditional on the source/measurement regime and frozen forecaster,
access rule, reserve fits and offsets. For a generic day under that
working law, the population counterparts are

\[
\begin{aligned}
c_m(\mathcal D)&=\frac{\mathbb E_{\mathcal D}[K^{-1}\sum_s A_{s,m,d}]}{\mathbb E_{\mathcal D}[N_d]},\\
v_{s,m}(\mathcal D)&=\frac{\mathbb E_{\mathcal D}[B_{s,m,d}]}{\mathbb E_{\mathcal D}[N_d]},\\
r_{m,b}(\mathcal D)&=\frac{c_m(\mathcal D)}{c_b(\mathcal D)}-1.
\end{aligned}
\]

The required expectations must be finite, expected eligible count must
be positive, and every population cost used as a ratio denominator must
be positive. Ratios of expected daily totals preserve the original
weighting by recorded turbine intervals; they are not unweighted means
of daily ratios. Missing-target occurrence is part of this working
observation law, not an assertion that missingness is random or that
unrecorded operating periods have the same risk.

These definitions make explicit the population interpretation required
by the existing resampling calculation; they do not change its score,
draws, thresholds or classifications. No process law or nominal-coverage
result is newly validated. The approximation requires a sufficiently
stable daily law and sufficiently short dependence for block resampling.
Seasonal composition, operating changes and clustered missingness can
violate it. If that interpretation fails, the finite-record point
comparisons and the classifications returned by the registered procedure
remain reproducible, but the nominal population-coverage interpretation
is not established.

The stored inference uses 20,000 shared circular day-block resamples,
bootstrap seed 20261029, primary span three days and separate
one-/seven-day sensitivities. All turbines, instances, policies and
conditions share the same sampled weather blocks. Numerators and
valid-target denominators remain separate and are recombined within each
draw. This preserves contemporaneous dependence rather than treating
turbines or learning seeds as independent weather replicates.

The simultaneous envelope uses the bootstrap maximum absolute
standardised deviation over its registered coordinates. Its critical
value is the higher empirical 0.95 quantile of that maximum. Non-exact
coordinates with constant resampled values remain unbounded;
algebraically exact self-ratios are zero. All bounds are dependent-data
approximations conditional on the five fitted instances. The primary and
full-q secondary families each use alpha=.05 separately; the article
does not supply one joint 95\% event covering both families and all
three block spans. Within a secondary family, all five tolerances reuse
its same bounds rather than generating five independent confidence
procedures. No distribution-free finite-sample guarantee or new coverage
experiment is claimed.

\subsection{S3.2. Primary fixed-reference
confirmation}\label{s3.2.-primary-fixed-reference-confirmation}

Let \(m\) be a historically selected candidate and \(b\) the frozen
qualified reference. Primary confirmation requires: (i) all-instance
reference qualification historically and in confirmation; (ii)
\(\max_s\overline V_{s,m}\le0.10\); and (iii)
\(\overline{(C_m/C_b-1)}\le0.02\). Here overbars denote the
corresponding existing joint upper bounds. A missing historical
selection returns no selection, and an unsuccessful condition remains
unconfirmed. No confirmation-side candidate replaces the original
choice.

The primary joint family contains only P1/HGB/q0.95 and WPQ36/HGB/q0.95:
ten instance event rates and four paired mean-cost coordinates,
including the self-ratios, for 14 coordinates. The result is supported
noninferiority within the declared 2\% research tolerance. It does not
test cost improvement, every-instance cost noninferiority or proximity
to every member of the full q family.

The observed costs are 67.020403 and 66.644430 kWh-equivalent per valid
interval for P1 and the reference, respectively. The relative premium is
0.5641\%, with approximate joint bounds {[}0.3575\%, 0.7708\%{]}. Both
bounds are positive. The maximum instance event-rate upper bounds are
6.2025\% and 6.2004\%, respectively. The positive premium remains part
of the result even though the 2\% criterion is met.

The existing point-accounting ledger also records mean uncovered
shortfall of 1.772331383 versus 1.734445047 interval-equivalent units:
an increase of 2.184349\% for P1, with positive point differences in all
five instances. Mean procurement is 49.297089493 versus 49.299979974.
The small event-rate decrease, from 5.5333784\% to 5.5313392\%,
therefore does not indicate lower uncovered-shortfall severity. No
additional interval or significance test is assigned to these ledger
quantities.

\subsection{S3.3. Secondary full-family sufficiency and
exclusion}\label{s3.3.-secondary-full-family-sufficiency-and-exclusion}

The full-q secondary analysis retains all 225 pitch-withheld policies
and all 65 policies in each pressure condition. Each span has 74,500
joint event-rate and pairwise cost coordinates across these six groups.
Calibration candidates were fitted before unlock, but selection of a
capable subset using confirmation outcomes is retrospective. It is
separate from the preselected primary confirmation.

For a given group, a policy remains possibly feasible when none of its
instance event-rate lower bounds exceeds 0.10. A policy is
risk-qualified when all its corresponding upper bounds are at most 0.10.
A risk-qualified policy is sufficient at tolerance \(\epsilon\) when the
upper bound on \(C_m/C_j-1\) is at most \(\epsilon\) against every
possibly feasible registered competitor \(j\). Competitors with
uncertain qualification are retained; historical nonqualification does
not establish their true infeasibility. The minimum resource tuple among
the sufficient policies gives the smallest \emph{qualified} set under
the declared order.

Identifying a finite-family resource minimum requires an additional
step. Every registered policy strictly smaller than the selected tuple
must be excluded either by an instance event-rate lower bound above 0.10
or by a cost-ratio lower bound above \(\epsilon\) against a
risk-qualified witness. A smaller policy with neither witness remains
unresolved. Thus an empty sufficient set, an unresolved smaller
alternative and an identified finite-family minimum are distinct
outcomes.

At 2\% in the three-day secondary analysis, P1/HGB/q0.925 has cost
65.739232, maximum instance event-rate upper bound 9.7699\% and resource
tuple \((1,1,5)\). Its stored worst cost upper bound against possibly
feasible competitors is 0.9742\%, so the cost-sufficiency step also
passes. All 15 strictly smaller registered policies are excluded. Table
S4 prints those existing exclusion records. The two risk exclusions are
Common/q0.90 policies; the remaining 13 use cost witnesses. These 15
policies cover the two Common estimators at all five offsets and
persistence at all five offsets. They do not include every conceivable
constant rule, architecture or alternative smaller function. Section S6
supplies the candidate and witness qualifications together.

\Needspace{18\baselineskip}

\textbf{Table S4. All smaller-policy exclusion witnesses for the
secondary q0.925 result.} Cost entries are existing lower bounds on the
smaller policy's relative cost premium against the named risk-qualified
witness, expressed as percentages; they exceed the 2\% margin. Risk
entries report one stored instance lower bound exceeding 10\%; the
accompanying CSV retains every witnessing instance. This table describes
the registered three-day full-q secondary family, not the primary q0.95
comparison.

{\def\LTcaptype{none} % do not increment counter
\begin{longtable}[]{@{}
  >{\raggedright\arraybackslash}p{(\linewidth - 6\tabcolsep) * \real{0.2500}}
  >{\raggedright\arraybackslash}p{(\linewidth - 6\tabcolsep) * \real{0.2500}}
  >{\raggedright\arraybackslash}p{(\linewidth - 6\tabcolsep) * \real{0.2500}}
  >{\raggedleft\arraybackslash}p{(\linewidth - 6\tabcolsep) * \real{0.2500}}@{}}
\toprule\noalign{}
\begin{minipage}[b]{\linewidth}\raggedright
Strictly smaller policy
\end{minipage} & \begin{minipage}[b]{\linewidth}\raggedright
Exclusion
\end{minipage} & \begin{minipage}[b]{\linewidth}\raggedright
Witness
\end{minipage} & \begin{minipage}[b]{\linewidth}\raggedleft
Existing lower bound (\%)
\end{minipage} \\
\midrule\noalign{}
\endhead
\bottomrule\noalign{}
\endlastfoot
Common/HGB/q0.90 & Event rate & Instance 201 & 10.2438 \\
Common/HGB/q0.925 & Relative cost & PQ6/HGB/q0.925 & 14.5901 \\
Common/HGB/q0.95 & Relative cost & PQ6/HGB/q0.925 & 15.4395 \\
Common/HGB/q0.975 & Relative cost & PQ6/HGB/q0.925 & 27.2390 \\
Common/HGB/q0.99 & Relative cost & PQ36/HGB/q0.925 & 61.2414 \\
Common/MLP/q0.90 & Event rate & Instance 201 & 10.0104 \\
Common/MLP/q0.925 & Relative cost & PQ6/HGB/q0.925 & 15.0854 \\
Common/MLP/q0.95 & Relative cost & PQ6/HGB/q0.925 & 15.8418 \\
Common/MLP/q0.975 & Relative cost & PQ6/HGB/q0.925 & 27.1526 \\
Common/MLP/q0.99 & Relative cost & PQ36/HGB/q0.925 & 61.6180 \\
Persistence/q0.90 & Relative cost & PQ36/HGB/q0.925 & 25.3079 \\
Persistence/q0.925 & Relative cost & PQ36/HGB/q0.925 & 25.0733 \\
Persistence/q0.95 & Relative cost & PQ36/HGB/q0.925 & 29.1403 \\
Persistence/q0.975 & Relative cost & PQ36/HGB/q0.925 & 46.5349 \\
Persistence/q0.99 & Relative cost & PQ36/HGB/q0.925 & 85.5821 \\
\end{longtable}
}

The primary P1/HGB/q0.95 policy is absent from this secondary full-q 2\%
sufficient set. Its noninferiority to one frozen reference cannot be
substituted for q0.925's finite-family result. Conversely, q0.925's
secondary identification cannot be described as confirmation that the
historical selector chose the full-family minimum. The registered
resource ordering also leaves sensor acquisition, network traffic, model
parameters, FLOPs and end-to-end execution outside the minimality claim.

The exclusion has a conditional algebraic interpretation. Suppose all
risk and relative-cost bounds used in a family cover their population
counterparts on one event \(E\). Let \(G\) be the truly risk-feasible
registered set, \(A\) the possibly feasible set and \(Q\) the
risk-qualified set; on \(E\), \(Q\subseteq G\subseteq A\). A candidate
qualified against every member of \(A\) is within the tolerance of the
best member of \(G\). A smaller candidate with a risk lower bound above
the cap is infeasible. If instead a smaller candidate \(j\) has a
cost-ratio lower bound above \(\epsilon\) against \(k\in Q\), then
\(c_j>(1+\epsilon)c_k\geq(1+\epsilon)\min_{g\in G}c_g\), so it cannot be
\(\epsilon\)-near-optimal. These implications explain why a richer
qualified witness can exclude a smaller candidate without circularity.
They do not establish the probability of \(E\) for the actual yearly
process.

\section{S4. Complete case records and negative
outcomes}\label{s4.-complete-case-records-and-negative-outcomes}

\subsection{S4.1. Hill of Towie: 90 dependent secondary
assessments}\label{s4.1.-hill-of-towie-90-dependent-secondary-assessments}

The complete secondary matrix has six information conditions, five
relative-cost tolerances (0\%, 0.5\%, 1\%, 2\% and 5\%) and three block
spans. Its 90 cases reuse the same observations, fitted instances and
policy predictions. The counts summarise classifications, not 90
independent experiments or a probability of deployment success.

Across these cases, 46 identify finite-family minima, 26 retain
unresolved smaller alternatives and 18 have no certified sufficient
candidate. The parallel assessment of the frozen historical sets gives
63 complete confirmations, 24 cases without a historical selection and
three unconfirmed cases. All three unconfirmed cases concern Delay 3 at
the 0.5\% tolerance, once for each block span. They fail the cost
criterion and are three sensitivity assessments of one condition, not
three independent failures. The secondary full-q q0.925 result does not
overwrite the frozen q0.95 failure at that tolerance.

The byte-identical attachment
\texttt{supplement\_data/towie\_all\_90\_cases.csv} retains the
condition, tolerance, span, sufficient-resource set,
smaller/excluded/unresolved counts and historical fixed-set status for
every case. In this file, \texttt{FINITE\_FAMILY\_MINIMUM\_IDENTIFIED}
means all required smaller candidates were excluded under the stored
bounds; \texttt{MINIMUM\_UNRESOLVED} means a sufficient set exists but
at least one smaller candidate was not excluded; and
\texttt{NO\_CERTIFIED\_SUFFICIENT\_POLICY} means no candidate met the
declared joint sufficiency rule. The last status does not prove that
every policy is truly infeasible.

\subsection{S4.2. WTB: 420 dependent development
assessments}\label{s4.2.-wtb-420-dependent-development-assessments}

WTB has 14 executed condition/budget settings: four pitch-withheld
budgets and two budgets for each of five pressure conditions. Crossing
these settings with the two registered analyses (historically selected q
recipes and the full-q grid), five tolerances and three spans produces
420 assessments. The 10,000-/40,000-row pressure settings were not
registered; their absence is not a failed experiment.

For the primary three-day span and 2\% tolerance, both registered
analyses across the 14 settings give 28 assessments. Every assessment
contains risk-qualified policies, but none establishes the required
near-optimal sufficiency. Across all 420 assessments, 367 have no
certified sufficient policy, 53 retain unresolved smaller alternatives
and none identifies a finite-family minimum. These outcomes concern
identification under a finite fitted family and approximate joint
bounds, not proof that more complex models or all feasible decision
rules are necessary.

At the prespecified 5\% sensitivity, the historical-q analysis
identifies a sufficient policy using current active and reactive power
with HGB and q=0.925 at all four pitch-withheld row budgets. At 360,000
rows, its mean cost is 32.670107 and its maximum instance risk upper
bound is 8.5952\%; 14 of 21 strictly smaller candidates remain
unresolved. At the same budget, P1 with HGB and q=0.925 costs 33.637629.
Its premium against PQ1 is 2.9615\%, with approximate joint bounds
{[}−2.2175\%, 8.1405\%{]}. These bounds neither qualify P1 within 5\%
nor exclude it as exceeding 5\%. The smallest currently qualified input
set is therefore not an identified minimum.

Two byte-identical attachments preserve this distinction: the
\href{supplement_data/wtb_all_420_sufficiency_cases.csv}{WTB sufficiency
cases} list every sufficient-set assessment, including the number of
risk-qualified policies; the
\href{supplement_data/wtb_all_420_minimum_cases.csv}{WTB minimum audits}
list the corresponding smaller-policy exclusions and unresolved labels.
They can be joined on condition, budget, span, analysis and tolerance.
Their rows describe the same dependent cases and must not be added
together as 840 experiments.

\Needspace{20\baselineskip}

\section{S5. Sensitivity, freezing and reproducibility
scope}\label{s5.-sensitivity-freezing-and-reproducibility-scope}

\subsection{S5.1. Separate block-length
sensitivities}\label{s5.1.-separate-block-length-sensitivities}

\textbf{Table S5. Stored full-q secondary identification under pitch
withholding at 2\%.} Cost is kWh-equivalent per valid turbine interval,
averaged over five fixed instances; event-rate upper bound is the
maximum across those instances. Each row uses a separate
74,500-coordinate approximate joint family. The rows do not form an
additional joint coverage statement.

{\def\LTcaptype{none} % do not increment counter
\begin{longtable}[]{@{}
  >{\raggedleft\arraybackslash}p{(\linewidth - 8\tabcolsep) * \real{0.2000}}
  >{\raggedright\arraybackslash}p{(\linewidth - 8\tabcolsep) * \real{0.2000}}
  >{\raggedleft\arraybackslash}p{(\linewidth - 8\tabcolsep) * \real{0.2000}}
  >{\raggedleft\arraybackslash}p{(\linewidth - 8\tabcolsep) * \real{0.2000}}
  >{\raggedleft\arraybackslash}p{(\linewidth - 8\tabcolsep) * \real{0.2000}}@{}}
\toprule\noalign{}
\begin{minipage}[b]{\linewidth}\raggedleft
Block span (days)
\end{minipage} & \begin{minipage}[b]{\linewidth}\raggedright
Identified policy
\end{minipage} & \begin{minipage}[b]{\linewidth}\raggedleft
Mean cost
\end{minipage} & \begin{minipage}[b]{\linewidth}\raggedleft
Maximum event-rate upper (\%)
\end{minipage} & \begin{minipage}[b]{\linewidth}\raggedleft
Smaller policies excluded
\end{minipage} \\
\midrule\noalign{}
\endhead
\bottomrule\noalign{}
\endlastfoot
1 & P1/HGB/q0.925 & 65.739232 & 9.6015 & 15/15 \\
3 & P1/HGB/q0.925 & 65.739232 & 9.7699 & 15/15 \\
7 & P1/HGB/q0.925 & 65.739232 & 9.8752 & 15/15 \\
\end{longtable}
}

The resource minimum is the same across these stored sensitivities, but
the risk upper bound approaches the 10\% research limit in the seven-day
calculation. The unchanged point cost follows from using the same
observations and policy; block length changes the uncertainty
calculation. Agreement across three spans does not establish nominal
coverage under arbitrary annual drift.

The pitch-withheld 5\% historical sets retain HGB and MLP ties, both
confirmed under all three spans. This sensitivity is not evidence that
neural estimators are universally inferior or redundant. Likewise, the
WTB and Hill of Towie differences do not isolate the effects of sample
count, terrain, control policy or a physical mechanism: several design
and source properties change together.

\subsection{S5.2. Actual pre-unlock
sequence}\label{s5.2.-actual-pre-unlock-sequence}

The completed executable/protocol was committed as \texttt{24153f8}
before confirmation. All 550 historical task bundles were verified,
including 3,240 archived members, common sample/target pairing and the
actual choices/freeze package. On 19 October 2026, the historical
package was created at 04:27:16.478898 UTC. Complete historical-result
review ran from 04:27:50 to 04:28:26, and actual freeze validation from
04:29:25 to 04:29:32. These checks preceded the exclusive unlock at
04:30:57.719354 UTC. The event binds the freeze used to open and process
the 2025 archive.

All 550 confirmation task bundles and their 2,200 members were checked
against stored manifests. The recovered compact result package contains
5,037 individually checksum-verified members. Read-only reviews
reproduced the main point quantities and the 90 secondary
classifications from existing stored records. No post-unlock refit, new
q level, threshold change, new evaluation window or new bootstrap draw
was used to obtain the reported results. These checks establish
provenance and stated consistency, not statistical coverage or field
safety.

The 2025 year was independent of the original frozen selector question
at unlock. Subsequent inspection for different explanations does not
retrospectively change that chronological fact. It does prevent
describing the same year as untouched confirmation of hypotheses
formulated after that inspection.

\subsection{S5.3. Available evidence and execution
boundaries}\label{s5.3.-available-evidence-and-execution-boundaries}

The original supplementary CSV attachments and their projection receipt
are inherited unchanged from the first version. They are literal
projections or byte-identical copies of completed outputs; the receipt
records input/output identities and row counts. V2's new Figure 3 fields
have a separate source manifest. None of these projections loads a
fitted model, predicts, calibrates an offset, resamples data or reads
new measurements. Table S6 identifies the portable package records for
checking the manuscript without rerunning the research analysis.

\Needspace{20\baselineskip}

\textbf{Table S6. Evidence locations in the portable review package.}
File areas are relative to that package's root. Public projections
explicitly omit internal deployment fields; their original identities
and transformations are recorded in PROVENANCE. These records are not a
full original execution archive.

{\def\LTcaptype{none} % do not increment counter
\begin{longtable}[]{@{}
  >{\raggedright\arraybackslash}p{(\linewidth - 2\tabcolsep) * \real{0.5000}}
  >{\raggedright\arraybackslash}p{(\linewidth - 2\tabcolsep) * \real{0.5000}}@{}}
\toprule\noalign{}
\begin{minipage}[b]{\linewidth}\raggedright
Evidence
\end{minipage} & \begin{minipage}[b]{\linewidth}\raggedright
Package area and file
\end{minipage} \\
\midrule\noalign{}
\endhead
\bottomrule\noalign{}
\endlastfoot
Execution settings and historical windows & \texttt{protocol/};
\texttt{confirmation\_protocol.json} (public projection) \\
Frozen choices and historical bounds & \texttt{data/history/};
\texttt{choices.json}, \texttt{bounds.json} \\
Primary labels, bounds and selection identity &
\texttt{data/confirmation/}; \texttt{primary.json} \\
Secondary full families and exclusions & Same area;
\texttt{secondary\_span1.json}, \texttt{secondary\_span3.json},
\texttt{secondary\_span7.json} \\
Freeze and unlock identity & \texttt{protocol/};
\texttt{freeze\_public\_projection.json},
\texttt{unlock\_public\_projection.json} \\
Procurement/shortage ledger & \texttt{data/ledger/};
\texttt{loss\_ledger.csv}, \texttt{policy\_seed\_ledger.csv} \\
WTB stored bounds and minimum audit & \texttt{data/wtb/};
\texttt{aggregate.json}, \texttt{necessity.json}; complete case tables
in \texttt{csv/} \\
Figure and field correspondence & \texttt{figure\_data/v2/}; four CSVs
and \texttt{fig3\_field\_manifest.json} \\
\end{longtable}
}

The confirmation receipts record CPU execution on Metis with Python
3.11.16, NumPy 2.4.6, pandas 2.3.3, scikit-learn 1.9.1 and PyTorch
2.11.0+cu128; the recorded PyTorch CUDA build does not imply GPU
training. Each task requested two CPUs, and its command, node, seed
settings and timestamps are retained in the execution receipts.
Requesting resources is not a measurement of actual resource
consumption.

Raw annual archives, complete physical arrays, full action arrays, fit
indices and model weights were retained on Metis; the compact local
recovery is not a second complete copy of those assets. The public
source route and exact code/configuration bindings support further
reconstruction, but a weight manifest is not proof that all original
weight bytes are currently available locally. Older graph, recovery and
mixture-of-experts branches have separate provenance gaps and are not
used to establish this article's primary comparison.

Finally, this task assumes synthetic release/arrival clocks and an
unchanged earlier forecast. It does not reconstruct measured
controller-message arrival times, point-of-common-coupling reserve
procurement, physical deliverability or an operator-approved service
level. The 2\% tolerance and 10\% event-rate cap are registered research
specifications. Reducing selected inputs from 112 to 5 does not
establish that sensors, upstream wind/power histories, SCADA
communication or actual costs can be removed. These boundaries define
what the observed trade-off can inform without converting it into a
deployment recommendation.

\Needspace{18\baselineskip}

\section{S6. Stored cost-sufficiency and witness
bridge}\label{s6.-stored-cost-sufficiency-and-witness-bridge}

\textbf{Table S7. Cost qualification, risk qualification and witness
status in the three-day full-q family.} Bounds are literal existing
fields expressed as percentages. A dash denotes a quantity not assigned
that role in this table, not zero uncertainty. Candidate rows use the
maximum paired-cost upper bound over all possibly feasible competitors.
The other rows identify the controlling competitor or qualified
exclusion witnesses. These are 74,500-coordinate secondary results,
separate from the primary 14-coordinate result.

{\def\LTcaptype{none} % do not increment counter
\begin{longtable}[]{@{}
  >{\raggedright\arraybackslash}p{(\linewidth - 6\tabcolsep) * \real{0.2500}}
  >{\raggedleft\arraybackslash}p{(\linewidth - 6\tabcolsep) * \real{0.2500}}
  >{\raggedleft\arraybackslash}p{(\linewidth - 6\tabcolsep) * \real{0.2500}}
  >{\raggedright\arraybackslash}p{(\linewidth - 6\tabcolsep) * \real{0.2500}}@{}}
\toprule\noalign{}
\begin{minipage}[b]{\linewidth}\raggedright
Role and policy
\end{minipage} & \begin{minipage}[b]{\linewidth}\raggedleft
Worst cost upper (\%)
\end{minipage} & \begin{minipage}[b]{\linewidth}\raggedleft
Maximum fit risk upper (\%)
\end{minipage} & \begin{minipage}[b]{\linewidth}\raggedright
Stored status
\end{minipage} \\
\midrule\noalign{}
\endhead
\bottomrule\noalign{}
\endlastfoot
Secondary candidate: P1/HGB/q0.925 & 0.9742 & 9.7699 & Sufficient;
smaller candidates excluded \\
Historical primary choice: P1/HGB/q0.95 & 3.5178 & 6.7049 &
Risk-qualified; cost not certified at 2\% \\
Controlling competitor: PQ36/HGB/q0.90 & --- & 11.5257 & Possibly
feasible, not risk-qualified \\
Cost-exclusion witness: PQ6/HGB/q0.925 & --- & 9.7564 &
Risk-qualified \\
Cost-exclusion witness: PQ36/HGB/q0.925 & --- & 9.6598 &
Risk-qualified \\
\end{longtable}
}

The two candidate cost upper bounds are controlled by PQ36/HGB/q0.90,
index 180. Its remaining possible feasibility is not proof of
feasibility. For the q0.95 candidate, the cost upper bounds against
already qualified PQ6/HGB/q0.925, PQ36/HGB/q0.925 and WPQ36/HGB/q0.925
are 3.4439\%, 3.4753\% and 3.4482\%, respectively. Their corresponding
lower bounds are 1.5795\%, 1.5651\% and 1.5235\%, all below 2\%. The
available evidence therefore fails to certify q0.95's full-family
sufficiency; it does not establish a true cost excess beyond 2\%.

Historical selection used a different stage and responsibility. In the
2024 selection record, P1/HGB/q0.925 had lower mean cost (72.047988)
than q0.95 (73.056204), but its fit-specific risk upper bounds ranged
from 11.0174\% to 11.0719\%, all above the qualification cap. The q0.95
upper bounds ranged from 7.8549\% to 7.8655\%, and it qualified
historically. These are upper-bound decisions, not proof that q0.925's
historical true risk exceeded 10\%. Its later secondary qualification
does not retrospectively change that selection.

The V2 figure-source CSV and field manifest retain exact policy indices
and paths for these stored values. The companion review package includes
the complete aggregates and qualification logic. No fitted model, new
draw or new confidence interval is needed to check this bridge.

\Needspace{15\baselineskip}

\section{S7. Prior knowledge and interface
qualification}\label{s7.-prior-knowledge-and-interface-qualification}

\textbf{Table S8. Actual overlap with close precedents.} Roles describe
the source's relation to this article; they do not imply that a prior
method fails. The remaining empirical difference is bounded by the
supplied forecast, fitted family and decision criteria.

{\def\LTcaptype{none} % do not increment counter
\begin{longtable}[]{@{}
  >{\raggedright\arraybackslash}p{(\linewidth - 4\tabcolsep) * \real{0.3333}}
  >{\raggedright\arraybackslash}p{(\linewidth - 4\tabcolsep) * \real{0.3333}}
  >{\raggedright\arraybackslash}p{(\linewidth - 4\tabcolsep) * \real{0.3333}}@{}}
\toprule\noalign{}
\begin{minipage}[b]{\linewidth}\raggedright
Source and role
\end{minipage} & \begin{minipage}[b]{\linewidth}\raggedright
Existing knowledge relevant here
\end{minipage} & \begin{minipage}[b]{\linewidth}\raggedright
Difference and limit of this article
\end{minipage} \\
\midrule\noalign{}
\endhead
\bottomrule\noalign{}
\endlastfoot
Pritchard 2011; mainstream & Current-power quantiles, extra covariates,
complete-year holdouts and different predictive attributes & A supplied
fixed forecast and a particular shortfall loss/risk task; no first claim
for power conditioning or cross-year input comparisons \\
Phipps et al.~2022; adjacent & Comparisons of post-processing position,
including a simpler final-power stage & Additional reserve-layer input
recipes and explicit qualification; not the first simpler downstream
procedure \\
Sun et al.~2022; mainstream & Point-forecast-derived reserves, coverage
and exceedance amounts, system consequences & Incremental inputs and
fixed-reference versus full-family responsibilities; no first claim for
reserve linkage or severity reporting \\
Serrano et al.~2024; adjacent & Feature choice by validation newsvendor
cost, bilevel method and synthetic experiments & Frozen fitted-family
qualification and exclusion, not a new feature-selection algorithm \\
Kowal 2022; foundational & Acceptable near-optimal families and smallest
acceptable subsets & Different conditional risk/cost bounds; no first
claim for an acceptable-set principle or new coverage theory \\
Dantas and Browell 2026; adjacent & Uncertainty-source comparisons
across horizons and test years & A fixed supplied commitment and
additional decision-layer inputs, not a new probabilistic forecasting
method \\
\end{longtable}
}

These comparisons are grounded in the original
papers\textsuperscript{3--8}. They do not certify exhaustive novelty.
The empirical addition is the observed cost/shortfall trade-off and the
different historical and retrospective qualifications in this controlled
interface.

The interface-material review separates three questions. First,
forecast-derived reserve layers and staged forecasting/dispatch are
documented, including Sun and colleagues' use of supplied point
forecasts. Second, that general practice does not establish this exact
one-forecast access contract and its \(t-6/t+6\) timing. The AEMO
overview, version 4 of August 2023, documents SCADA inputs, AWEFS/ASEFS
forecasts and NEMDE use, with five-minute forecast updates in the
dispatch/5MPD setting\textsuperscript{9}. It provides a different staged
interface rather than evidence for retaining this particular old
commitment; the document is not used to assert current 2026 operating
rules. Third, the material examined does not establish a realised
benefit of reducing the learner's raw columns from 112 to five,
particularly when the current implementation still constructs 148
columns and retains upstream W/P. The last two questions remain
unclosed.

The precise missing materials are an actual access/update specification
for the downstream service, timestamped message routes if real arrival
is claimed, a documented channel-priority or integration burden, and
evidence of an acceptable cost/severity exchange if deployment benefit
is claimed. The current paper does not claim these benefits. Operator
documentation, interface schemas, owner-provided requirements or a later
separately authorised implementation assessment would be appropriate
acquisition routes; the experimental protocol cannot serve as external
proof of those operational conditions.

\section{S8. Scope of the portable review
package}\label{s8.-scope-of-the-portable-review-package}

The package supports offline verification of completed aggregates,
candidate identities, classification inequalities, cost/shortfall
accounting and figure-source values. It does not train, run a fitted
model, calibrate offsets, create bootstrap draws or read new
measurements. Complete stored statistical aggregates are retained where
safe; deployment-specific paths or host fields are omitted only in
explicitly labelled public projections with original source identities.
Such a projection is not presented as a byte-identical raw receipt.

The package is prepared for local review before an author-approved
public release. No permanent public URL or new licence grant has been
fabricated. Original annual measurements, full action arrays, fit
indices, weights, copyright paper PDFs and private skill files are not
distributed. The absence of those assets bounds the claim to
stored-result verification rather than independent re-estimation of the
entire study.

\section*{References}\label{references}
\addcontentsline{toc}{section}{References}

\protect\phantomsection\label{refs}
\begin{CSLReferences}{0}{1}
\bibitem[\citeproctext]{ref-clerc2026towie}
\CSLLeftMargin{1. }%
\CSLRightInline{Clerc A, Lingkan E. Hill of towie wind farm open
dataset. Published online 2026.
doi:\href{https://doi.org/10.5281/zenodo.22662930}{10.5281/zenodo.22662930}}

\bibitem[\citeproctext]{ref-zhou2024sdwpfdataset}
\CSLLeftMargin{2. }%
\CSLRightInline{Zhou J. {SDWPF}: A dataset for spatial dynamic wind
power forecasting over a large turbine array. Published online 2024.
doi:\href{https://doi.org/10.6084/m9.figshare.24798654.v2}{10.6084/m9.figshare.24798654.v2}}

\bibitem[\citeproctext]{ref-pritchard2011shortterm}
\CSLLeftMargin{3. }%
\CSLRightInline{Pritchard G. Short-term variations in wind power: Some
quantile-type models for probabilistic forecasting. \emph{Wind Energy}.
2011;14(2):255-269.
doi:\href{https://doi.org/10.1002/we.416}{10.1002/we.416}}

\bibitem[\citeproctext]{ref-phipps2022postprocessing}
\CSLLeftMargin{4. }%
\CSLRightInline{Phipps K, Lerch S, Andersson M, Mikut R, Hagenmeyer V,
Ludwig N. Evaluating ensemble post-processing for wind power forecasts.
\emph{Wind Energy}. 2022;25(8):1379-1405.
doi:\href{https://doi.org/10.1002/we.2736}{10.1002/we.2736}}

\bibitem[\citeproctext]{ref-sun2022dynamicreserve}
\CSLLeftMargin{5. }%
\CSLRightInline{Sun Y, Nelson JH, Stevens JC, et al. Machine learning
derived dynamic operating reserve requirements in high-renewable power
systems. \emph{Journal of Renewable and Sustainable Energy}.
2022;14(3):036303.
doi:\href{https://doi.org/10.1063/5.0087144}{10.1063/5.0087144}}

\bibitem[\citeproctext]{ref-serrano2024featureselection}
\CSLLeftMargin{6. }%
\CSLRightInline{Serrano B, Minner S, Schiffer M, Vidal T. Bilevel
optimization for feature selection in the data-driven newsvendor
problem. \emph{European Journal of Operational Research}.
2024;315(2):703-714.
doi:\href{https://doi.org/10.1016/j.ejor.2024.01.025}{10.1016/j.ejor.2024.01.025}}

\bibitem[\citeproctext]{ref-kowal2022subset}
\CSLLeftMargin{7. }%
\CSLRightInline{Kowal DR. Bayesian subset selection and variable
importance for interpretable prediction and classification.
\emph{Journal of Machine Learning Research}. 2022;23(108):1-38.
\url{https://www.jmlr.org/papers/v23/21-0403.html}}

\bibitem[\citeproctext]{ref-dantas2026seamless}
\CSLLeftMargin{8. }%
\CSLRightInline{Dantas G, Browell J. Seamless short- to mid-term
probabilistic wind power forecasting. \emph{Wind Energy}.
2026;29(2):e70079.
doi:\href{https://doi.org/10.1002/we.70079}{10.1002/we.70079}}

\bibitem[\citeproctext]{ref-aemo2023awefs}
\CSLLeftMargin{9. }%
\CSLRightInline{Australian Energy Market Operator. \emph{Australian Wind
Energy Forecasting System ({AWEFS}) and Australian Solar Energy
Forecasting System ({ASEFS})}. Australian Energy Market Operator; 2023.
\url{https://aemo.com.au/-/media/files/electricity/nem/security_and_reliability/dispatch/policy_and_process/australian-wind-energy-forecasting-solar-energy-forecasting-system.pdf}}

\end{CSLReferences}

\end{document}
